% ============================================================
% VI. CONCLUSIONES
% ============================================================

\section{Conclusiones}

La Transformada Rápida de Fourier, calculada mediante \texttt{signal\_lab.spectral.compute\_fft}, permitió transformar las señales de audio del dominio temporal al dominio frecuencial, facilitando la identificación de componentes espectrales característicos de cada muestra dentro de la ventana de análisis de 3 s (o de la grabación completa, cuando esta es más corta).

En los sonidos animales, procesados a partir de grabaciones reales de dominio público/CC, se observaron diferencias claras y cuantificables entre las frecuencias dominantes del gato ($1545.17$ Hz), el perro ($459.73$ Hz) y el gallo ($1731.67$ Hz): los tres pares superaron ampliamente el umbral de separabilidad clara ($S_{a,b} \geq 3$) definido en la Ec.~\ref{eq:separabilidad}, con puntajes entre $442$ y $4267$. Para este conjunto, tanto la lectura cualitativa (Fig.~\ref{fig:animales_comparacion}) como la cuantitativa coinciden en que las tres señales son distinguibles por su comportamiento espectral.

Los instrumentos musicales y las voces humanas, en cambio, se mantienen en este informe como valores de referencia del borrador original, pendientes de grabaciones reales; sobre esos valores se observó una tendencia creciente en la frecuencia dominante desde el contrabajo hasta la flauta, y diferencias espectrales entre las dos voces pese a compartir la misma frase, pero estas lecturas deben confirmarse con datos reales antes de considerarse comparables a los resultados de animales.

La media y la desviación estándar, calculadas de forma aislada, no resultaron características suficientes para identificar las diferentes señales; solo cobran utilidad como parte del denominador del puntaje de separabilidad. Precisamente al recalcular ese puntaje con datos reales se identificó una debilidad de diseño (Sección~\ref{sec:discusion-unidades}): al dividir una diferencia de frecuencia (Hz) por una suma de desviaciones estándar de amplitud (adimensional), su magnitud depende fuertemente de la escala de preprocesamiento de la señal, por lo que sus umbrales fijos solo son comparables entre señales procesadas por el mismo pipeline de normalización.

Finalmente, los resultados de instrumentos y voces siguen siendo exploratorios por no contar aún con grabaciones propias, y los de animales, aunque ya verificados con datos reales, provienen de una sola grabación por categoría. Los próximos pasos naturales son: (1) sustituir los valores de referencia de instrumentos y voces siguiendo el mismo procedimiento usado para los animales (descargar o grabar audio, documentarlo en \texttt{audio\_samples/SOURCES.md}, y ejecutar el pipeline \texttt{load\_audio} + \texttt{analyze\_audio}); (2) recolectar varias grabaciones independientes por categoría; (3) corregir la Ec.~\ref{eq:separabilidad} para que ambos términos compartan unidades comparables; y (4) incorporar características espectrales adicionales (armónicos, energía por bandas) antes de evaluar un clasificador automático de fuentes sonoras.
